\section{Conclusions}\label{sec:conclusion}

The search for hidden attractors has been formulated as the detection of a rare event in the space of initial conditions, with the detection probability $P(N)$ and the cost $N_{0.9}$ relative to uniform sampling as the measures of quality. We have built a four-level benchmark with an identification and verification protocol and synthetic families with exactly known basin measure, and we have proposed a portfolio of multilevel splitting driven by transient duration, active learning of basin boundaries, and a uniform share, which in the worst case is at most $1/\varepsilon$ times costlier than uniform sampling and inherits the scaling exponent of the splitting arm.

On rare hidden attractors with $p\sim10^{-3}$--$10^{-2}$, the portfolio reduces the detection cost by factors of 3--7, with $R=100$ runs per problem. On a synthetic family with a graded signal, the cost grows as $p^{-0.45}$ instead of $p^{-1}$ down to $p=10^{-5}$; below this, the time resolution of the reaction coordinate limits the gain, and a proportionally smaller integration step restores it. On coupled copies of systems with hidden attractors, the gain in integrations reaches $10^3$--$10^4$. Counted in wall time, the gain depends on how expensive the nontarget trajectories are: where trajectories are bounded, as in the Chua circuits, it carries over, and on two coupled Chua circuits a rare hidden regime ($p=2\times10^{-5}$) is found 64 times faster than by uniform sampling with the same evaluation and 10 times faster than by uniform sampling with recurrence-based early stopping; where most trajectories escape within a few steps, a uniform point is cheap, and the gain in time is much smaller than in integrations or absent. The enumeration of regimes by the identification protocol revealed a hidden chaotic attractor of the coupled Chua circuits that is born by the coupling at $\epsilon\approx0.045$ and persists to $\epsilon\approx0.5$. Enlarging the region of initial conditions and continuation in a parameter, in contrast, do not make the hidden attractors of the benchmark rare. Failure cases are documented explicitly, and in all of them the portfolio stays within the guarantee.
